% =========================================================================
% Figura (apuntes manuscritos, pág. 5, y de teoría, pág. II-4): región V'
% polarizada; el elemento dV' en r' contiene el dipolo dp = P dV', que crea
% en P (posición r) el potencial dphi = dp.u_R/(4 pi eps0 R^2).
% =========================================================================
\begin{figure}[ht]
\centering
\begin{tikzpicture}[x=1cm, y=1cm, >=Stealth]
  \ejesO{4.4}{6.6}{1.6}
  \regionV
  % pequeños dipolos (polarización) en el interior
  \foreach \px/\py in {1.45/2.5, 1.75/3.35, 2.55/3.55, 3.15/3.1, 3.2/2.25, 2.45/1.95, 1.75/1.95}{%
    \draw[-{Stealth[length=2.6pt,width=2.2pt]}, magacero!75, line width=0.55pt] (\px-0.12,\py-0.08) -- (\px+0.14,\py+0.1);}
  \node[etiq, text=magprof] at (0.8,3.65) {$V'$};
  % elemento de volumen
  \coordinate (Q) at (2.35,2.75);
  \coordinate (P) at (6.2,3.25);
  \draw[vect] (0,0) -- ($(Q)+(-0.1,-0.12)$) node[pos=0.5, above left=-1pt, etiq] {$\rp$};
  \cubito{Q}
  \node[etiq, font=\footnotesize, anchor=south east] at ($(Q)+(-0.1,0.12)$) {$\dVp$};
  \draw[vect] (0,0) -- (P) node[pos=0.72, below right=-2pt, etiq] {$\vc{r}$};
  \draw[vectrojo] ($(Q)+(0.22,0.06)$) -- (P) node[pos=0.6, above=2pt, etiq, text=rojoteo] {$\vc{R}=\vc{r}-\rp$};
  \fill[magprof] (P) circle (1.5pt) node[right=2pt, etiq] {$P$};
  % normal exterior
  \coordinate (B) at (3.3,1.8);
  \draw[vectverde] (B) -- ++(0.42,-0.44) node[below right=-2pt, etiq, text=verdeteo] {$\un$};
  \node[etiq, font=\footnotesize, text=magacero, anchor=south] at (2.45,3.98) {$S'=\partial V'$};
  \fill[black] (0,0) circle (1.2pt) node[below right=1pt, etiq, font=\footnotesize] {$O$};
  % recuadro con las relaciones
  \node[cuadro, anchor=north west, align=left, font=\footnotesize, inner sep=4pt] at (7.0,2.6)
    {$\dif\vp=\vP(\rp)\,\dVp$\\[4pt]
     $\dif\phi(\vc{r})=\dfrac{\dif\vp\cdot\uR}{4\pi\eo R^{2}}$\\[7pt]
     $\rhoP=-\nabp\cdot\vP$,\quad $\sigP=\vP\cdot\un$};
\end{tikzpicture}
\caption{Distribución continua de dipolos en una región $V'$ limitada por la superficie $S'$, de normal exterior $\un$. El elemento $\dVp$ situado en $\rp$
contiene el momento dipolar $\dif\vp=\vP\,\dVp$, que contribuye al potencial en el punto campo $P$ como un dipolo puntual a la distancia $R$. El potencial
total equivale al de las cargas de polarización $\rhoP$ y $\sigP$ (\cref{teo:cargas-pol}).}
\label{fig:polarizacion}
\end{figure}
